Measurements of charged lepton scattering
from different nuclei show that the cross section ratio on a heavy nucleus relative to 
the deuteron $\sigma^{\A}/\sigma^{D}$ deviates from unity by as much as 20\%. 
This demonstrates nontrivial nuclear effects over a wide range of Bjorken's scaling variable \Xbj~\cite{Arneodo:1992wf,geesaman1995nuclear,Norton:2003cb,rith2014}.
These observations, first reported by the European Muon Collaboration (EMC)~\cite{Aubert:1983xm,Aubert1987740} in 1983, signal a difference in the quark-parton structure of bound nucleons from that of free nucleons
and have triggered theoretical exploration of
background nuclear mechanisms~\cite{Arneodo:1992wf,Norton:2003cb}.

In neutrino physics, understanding nuclear effects is necessary for correct
interpretation of measurements of electroweak parameters 
and evaluation of corresponding uncertainties~\cite{PhysRevLett.88.091802}.
The precision of modern neutrino oscillation experiments
has rekindled interest in measuring nuclear effects, albeit at lower 
neutrino energies where elastic and resonance processes, rather than deep inelastic processes, are dominant~\cite{Coloma:2013tba}.

Neutrino scattering, unlike that of charged leptons, involves the axial-vector current and
is sensitive to specific quark and antiquark flavors.
Therefore, nuclear modifications of neutrino cross sections may differ
from those of charged leptons~\cite{Kulagin2006126,PhysRevC.84.054610,Haider:2013eqa,Kopeliovich:2012kw}.
An indirect extraction of neutrino deep inelastic structure function ratios using NuTeV Fe~\cite{PhysRevD.74.012008} and
CHORUS Pb~\cite{onegut200665} data suggests this is the case~\cite{Kovarik:2010uv}.
If confirmed, this either challenges the validity of QCD factorization for processes
involving bound nucleons or signals inconsistency between neutrino and charged lepton data.
Another study~\cite{PhysRevLett.110.212301} using different techniques does not find this behavior.
Neutrino scattering data are necessary for separation of valence and sea quark contributions to parton distribution functions~\cite{PhysRevC.76.065207,PhysRevD.77.054013,Kovarik:2010uv}, but high-statistics data from iron and lead must be corrected to account for poorly measured nuclear modifications.

Direct measurements of neutrino cross section ratios for different nuclei
are therefore of significant interest and importance.
So far, the only such measurements are ratios of Ne to D~\cite{Parker19841,Cooper1984133,PhysRevD.32.2441}, 
but these are rarely used because of large statistical uncertainties and model-dependent extraction from a mixed H-Ne target.
In this Letter, we report the first measurement of inclusive charged-current neutrino cross section
ratios of C, Fe, and Pb to scintillator (CH) as functions of neutrino energy \Enu and \Xbj.
This is the first application to neutrino physics of the EMC-style technique of measuring nuclear dependence with multiple nuclear targets in the same beam and detector.


